\section*{Capítulo \ref{ch:b3:water-rock-interaction} --- Interação água--rocha}

\begin{solution}{exo:b3:water-rock-interaction:1}
Cátions: $2\times48/40.08 + 2\times9.7/24.31 + 23/22.99 + 3.9/39.10 =
\qty{4.29}{meq/L}$; ânions: $183/61.02 + 35.5/35.45 + 2\times48/96.06 =
\qty{5.00}{meq/L}$. Erro $100\times(4.29 - 5.00)/9.29 = \qty{-7.6}{\percent}$:
muito acima do 1 ou \qty{2}{\percent} de uma análise cuidadosa de uma água desse tipo.
Falta um cátion ou está subestimado, ou um ânion está alto demais.
\end{solution}

\begin{solution}{exo:b3:water-rock-interaction:2}
$\sqrt I = 0.0707$. \ce{Ca^2+}: $-\log\gamma = 0.5085\times4\times0.0707/(1 +
0.3281\times6\times0.0707) = 0.126$, $\gamma = 0.748$; \ce{HCO3-}: $\gamma =
0.927$; \ce{SO4^2-}: $\gamma = 0.739$.
\end{solution}

\begin{solution}{exo:b3:water-rock-interaction:3}
$\sqrt I = 0.126$; $\gamma_{\ce{Ca}} = 0.622$, $\gamma_{\ce{SO4}} = 0.602$;
$\log\mathrm{IAP} = \log(3.0\times10^{-3}\times0.622\times4.0\times10^{-3}
\times0.602) = -5.35$; $\mathrm{SI} = -5.35 + 4.55 = -0.80$. A água está
subsaturada e ainda pode dissolver gipsita.
\end{solution}

\begin{solution}{exo:b3:water-rock-interaction:4}
$m = 0.012/60.08 = \qty{2.00e-4}{mol/L}$, $\log m = -3.70$. Quartzo:
$\mathrm{SI} = -3.70 + 3.98 = +0.28$, levemente supersaturada; sílica amorfa:
$-3.70 + 2.71 = -0.99$, subsaturada.
\end{solution}

\begin{solution}{exo:b3:water-rock-interaction:5}
A \qty{25}{\celsius}, $\log K = -8.45 - 6.35 - 1.47 + 10.33 = -5.94$, $m =
(10^{-5.94}\times10^{-2}/4)^{1/3} = \qty{1.42}{mmol/L}$, \qty{57}{mg/L}; a
\qty{10}{\celsius}, $\log K = -5.59$, $m = \qty{1.86}{mmol/L}$,
\qty{75}{mg/L}: um terço a mais no frio.
\end{solution}

\begin{solution}{exo:b3:water-rock-interaction:6}
$[\ce{CO2}] = 10^{-7.60}\times4.0\times10^{-3}/10^{-6.35} =
\qty{2.25e-4}{mol/L}$; $[\ce{CO3^2-}] = 10^{-10.33}\times4.0\times10^{-3}/
10^{-7.60} = \qty{7.4e-6}{mol/L}$; $\log P_{\ce{CO2}} = -3.648 + 1.47 =
-2.18$, 16 vezes o ar. Total \qty{4.23}{mmol/L}: \ce{CO2}
\qty{5.3}{\percent}, \ce{HCO3-} \qty{94.5}{\percent}, \ce{CO3^2-}
\qty{0.18}{\percent}.
\end{solution}

\begin{solution}{exo:b3:water-rock-interaction:7}
$r = 10^{-12.09} = \qty{8.2e-13}{mol.m^{-2}.s^{-1}}$ (os termos ácido, neutro e
básico contribuem $10^{-12.45}$, $10^{-12.56}$ e $10^{-12.74}$). O raio
diminui à taxa $rV_m = 8.2\times10^{-13}\times1.0131\times10^{-4} =
\qty{8.3e-17}{m/s}$, \qty{2.6}{nm} por ano; o raio de \qty{0.5}{mm} dura
$6.0\times10^{12}$~s, cerca de \num{190000} anos.
\end{solution}

\begin{solution}{exo:b3:water-rock-interaction:8}
$\exp[(69.8/0.008314)(1/278.15 - 1/298.15)] = 7.6$; com
\qty{74}{kJ/mol}, 8.6.
\end{solution}

\begin{solution}{exo:b3:water-rock-interaction:9}
Igualando $(12.703 - X - 3Y)/3$ e $(7.435 - 2Y)/2$ elimina-se $Y$: $X =
12.703 - 11.153 = 1.55$, uma linha horizontal. A água tem $X = \log(2.0\times
10^{-3}/39.10) + 7.0 = 2.71$ e $Y = \log(0.015/60.08) = -3.60$: acima de 1.55
e bem abaixo da linha muscovita--feldspato-K ($X = -3.161 - 3Y = 7.65$), logo
forma muscovita.
\end{solution}

\begin{solution}{exo:b3:water-rock-interaction:10}
$\delta_w^f = -2 - 2 = -4$. Fechado: $W/R = 0.5\times(-2 - 6)/(-13 + 4) =
0.44$. Aberto: $0.5\ln[(-8)/(-9) + 1] = 0.5\ln1.889 = 0.32$. No sistema aberto,
cada porção sai em equilíbrio com a rocha que encontrou por último, levando
consigo o máximo de oxigênio-18 possível; no sistema fechado, todo o lote
termina em $-4$ e retira menos por quilograma.
\end{solution}

\begin{solution}{exo:b3:water-rock-interaction:11}
Aberto: $\log K = -5.59$, $m = (10^{-5.59}\times10^{-1.5}/4)^{1/3} =
\qty{2.73}{mmol/L}$, \qty{109}{mg/L}. Fechado: $C_0 = 10^{-1.27 - 1.5} =
\qty{1.70e-3}{mol/L}$ e $K_{\mathrm{cal}}K_1/K_2 = 10^{-4.32}$; a cúbica
$4m^3 = 4.79\times10^{-5}(1.70\times10^{-3} - m)$ dá $m =
\qty{1.45}{mmol/L}$, \qty{58}{mg/L}, cerca da metade. O dióxido de carbono
restante é $0.25\times10^{-3}$, \qty{15}{\percent} de $C_0$, de modo que
$\log P_{\ce{CO2}}$ cai de $-1.5$ para $-2.33$.
\end{solution}

\begin{solution}{exo:b3:water-rock-interaction:12}
$1 - 10^{\mathrm{SI}}$: 0.99, 0.90 e 0.21. A lei de potência dá
$(10^4/6)^{-0.61} = 0.011$: uma taxa cerca de 90 vezes menor. Um mineral de
campo é antigo, com as superfícies reativas já consumidas, e está banhado por
água porosa lenta e próxima do equilíbrio; os dois fatores se multiplicam, e a
taxa de campo cai ordens de grandeza abaixo da taxa de laboratório.
\end{solution}

\begin{solution}{pb:b3:water-rock-interaction:1}
\textbf{1.} \ce{Ca} 2.795, \ce{Mg} 0.165, \ce{Na} 0.305, \ce{K} 0.031,
\ce{HCO3} 5.572, \ce{Cl} 0.310, \ce{SO4} 0.146, \ce{NO3} \qty{0.145}{mmol/L}.

\textbf{2.} Cátions \qty{6.25}{meq/L}, ânions \qty{6.32}{meq/L}: erro
\qty{-0.5}{\percent}, aceitável.

\textbf{3.} $I = \tfrac12(4\times2.795 + 4\times0.165 + 0.305 + 0.031 + 5.572
+ 0.310 + 4\times0.146 + 0.145)\times10^{-3} = 0.0094$.

\textbf{4.} $\sqrt I = 0.0969$: $\gamma_{\ce{Ca^2+}} = 0.683$,
$\gamma_{\ce{HCO3-}} = 0.904$, $\gamma_{\ce{CO3^2-}} = 0.676$.

\textbf{5.} Sim: $I$ é um décimo do limite de cerca de 0.1.

\textbf{6.} $a_{\ce{HCO3-}} = 5.572\times10^{-3}\times0.904 = 5.04\times
10^{-3}$ ($\log = -2.298$).

\textbf{7.} $a_{\ce{CO2}} = 10^{-7.20}\times5.04\times10^{-3}/10^{-6.46} =
9.17\times10^{-4}$ ($\log = -3.038$).

\textbf{8.} $\log P_{\ce{CO2}} = -3.038 + 1.27 = -1.77$, \qty{0.017}{atm}:
$10^{1.60} = 40$ vezes o ar, dentro da faixa das águas subterrâneas. O dióxido de carbono
provém do ar do solo da área de recarga, das raízes e da decomposição.

\textbf{9.} $a_{\ce{CO3^2-}} = 10^{-10.49}\times5.04\times10^{-3}/10^{-7.20}
= 2.58\times10^{-6}$, $[\ce{CO3^2-}] = \qty{3.8e-6}{mol/L}$. Total
\qty{6.49}{mmol/L}: \ce{CO2} \qty{14.1}{\percent}, \ce{HCO3-}
\qty{85.8}{\percent}, \ce{CO3^2-} \qty{0.06}{\percent}.

\textbf{10.} Alcalinidade $5.572 + 2\times0.0038 = \qty{5.580}{meq/L}$:
\qty{99.9}{\percent} de bicarbonato.

\textbf{11.} $a_{\ce{Ca^2+}} = 2.795\times10^{-3}\times0.683 = 1.909\times
10^{-3}$; $\log\mathrm{IAP} = \log(1.909\times10^{-3}\times2.58\times10^{-6})
= -8.307$; $\mathrm{SI} = -8.307 + 8.35 = +0.04$.

\textbf{12.} Dentro de $\pm0.1$, a incerteza do índice: a água está em
equilíbrio com a calcita. Ela permaneceu muito tempo na rocha: uma nascente
lenta e difusa.

\textbf{13.} $m = [10^{-5.59}\times0.0171/(4\times0.683\times0.904^2)]^{1/3}
= \qty{2.70}{mmol/L}$, \qty{108}{mg/L}, contra \qty{2.79}{mmol/L} medidos:
\qty{3}{\percent} a mais, como indica um índice de $+0.04$.

\textbf{14.} $a_{\ce{CO2}} = 10^{-1.27 - 3.37} = 2.29\times10^{-5}$;
$a_{\ce{H+}} = 10^{-6.46}\times2.29\times10^{-5}/5.04\times10^{-3}$, pH 8.80;
$a_{\ce{CO3^2-}} = 1.03\times10^{-4}$; $\mathrm{SI} = \log(1.909\times
10^{-3}\times1.03\times10^{-4}) + 8.35 = +1.65$ (1.61 com o balanço
completo de alcalinidade da figura).

\textbf{15.} Iterando $m = [10^{-5.59 - 3.37}/(4\gamma_{\ce{Ca}}
\gamma_{\ce{HCO3}}^2)]^{1/3}$ com $I = 3m$: $m = \qty{0.72}{mmol/L}$
(\qty{29}{mg/L}), $I = 0.0022$. A calcita precipitou: $2.795 - 0.718 =
\qty{2.08}{mmol/L}$, cerca de \qty{210}{mg} por litro.

\textbf{16.} $0.208\times2.0 = \qty{0.42}{g/s}$, \qty{13}{t} por ano,
$13.1/2.71 = \qty{4.8}{m^3}$; sobre \qty{400}{m^2}, \qty{12}{mm} por ano,
\qty{0.03}{mm} por dia: bem abaixo dos tufos de nascentes frias citados no capítulo.

\textbf{17.} $X = \log(1.2\times10^{-3}/39.10) + 5.8 = 1.29$; $Y =
\log(9.0\times10^{-3}/60.08) = -3.82$.

\textbf{18.} $X < 1.55$ e $Y > -4.39$: caulinita. Quartzo: $\mathrm{SI} =
-3.82 + 3.98 = +0.16$.

\textbf{19.} Abaixo de $\log a_{\ce{H4SiO4}} = -4.39$, \qty{2.4}{mg/L} de
sílica.

\textbf{20.} $X = 1.55$ requer pH $= 1.55 + 4.51 = 6.06$.

\textbf{21.} Ácido $-10.06 - 0.500\times5.8 = -12.96$; neutro $-12.41$; básico
$-21.20 + 0.823\times5.8 = -16.43$. Soma: $\log r = -12.30$, com o termo neutro
dominante.

\textbf{22.} \omterm{def:b3:water-rock-interaction:regimes}{Regime limitado pelo suprimento}: solo espesso e encostas suaves significam erosão lenta;
o feldspato fresco que alcança a frente de intemperismo é totalmente lixiviado, e
o fluxo acompanha a taxa de erosão, e não o clima.

\textbf{23.} $\exp[(74/0.008314)(1/283.15 - 1/298.15)] = 4.9$.

\textbf{24.} O índice de saturação da calcita na água da nascente em sua saída
é $\mathrm{SI} = \textbf{+0.04}$.
\end{solution}
